\documentclass[11pt,a4paper]{article}
\usepackage[T1]{fontenc}
\usepackage[utf8]{inputenc}
\usepackage{lmodern}
\usepackage[margin=27mm]{geometry}
\usepackage{amsmath,amssymb,amsthm,mathtools}
\usepackage{microtype}
\usepackage[numbers,sort&compress]{natbib}
\usepackage{xcolor}
\usepackage[colorlinks=true,linkcolor=blue!45!black,citecolor=blue!45!black,urlcolor=blue!45!black]{hyperref}
\newtheorem{theorem}{Theorem}
\newtheorem{lemma}{Lemma}
\newcommand{\K}{\mathcal K}
\newcommand{\spec}{\operatorname{spec}}
\newcommand{\Rea}{\operatorname{Re}}
\newcommand{\Ima}{\operatorname{Im}}
\newcommand{\av}[1]{\overline{#1}}
\newcommand{\ii}{\mathrm i}
\newcommand{\Bstar}{B_*}
\newcommand{\Nloc}{\mathcal N}
\newcommand{\trans}{^{\mathsf T}}
\setlength{\emergencystretch}{2em}
\title{A missing spectral arc at finite cycle affinity\protect\\ in four-state Markov generators}
\author{Xiaojian Zeng\thanks{Email: \href{mailto:xiaojianzeng@foxmail.com}{xiaojianzeng@foxmail.com}.}\\[3pt]
\small PBC School of Finance, Tsinghua University, 100190}
\date{}

\begin{document}
\maketitle

\begin{abstract}
Cycle affinities constrain the decay rates and oscillation frequencies of a
Markov process, but a spectral localization bound need not describe all
attainable eigenvalues. We disprove the proposed equality between the
finite-affinity spectral set and a vertical compression of the unrestricted
set, already for four states. At maximum allowed escape rate one and cycle
budget $B_*=(\log 3)/2$, an entire sufficiently small compressed Karpelevich
arc near $-1+\ii/2$ is absent. The exclusion holds in a compact class defined
by all oriented cycle-product inequalities, including reducible and
zero-rate limits. The proof combines the equality case of cycle-winding
localization with a uniform quadratic expansion about a biased square.
Opening chords creates triangles whose affinity constraints give an upper
quadratic coefficient $1/5$, whereas the compressed arc has coefficient
$1/2$. A coercive estimate controls unequal escape rates, unused affinity,
and arbitrary approaches to the endpoint. The validity interval is
existential: no numerical radius, sharpness of $1/5$, or full boundary
classification is asserted.
\end{abstract}

\noindent\textbf{Keywords:} Markov generators; cycle affinity; spectral
localization; nonequilibrium relaxation; inverse spectral problem.

\section{Spectral localization and attainability}

An eigenvalue of a continuous-time Markov generator specifies a possible
decay rate and, when it is nonreal, an oscillation frequency. Thermodynamic
restrictions on these quantities provide information about relaxation
without specifying all transition rates. Uhl and Seifert proposed an
affinity-dependent spectral ellipse and studied perturbations of a uniform
biased cycle \cite{UhlSeifert2019}. Xu \emph{et al.}\ established an ellipse
based on edge forces and posed a stronger question about the complete
spectral set: does vertical compression of the unrestricted region give
exactly the finite-force region \cite[End Matter and Supplementary
Eq.~(S3)]{Xu2025}? The later cycle-winding localization theorem of
Kolchinsky, Ohga and Ito proves the cycle-affinity ellipse for general
finite generators \cite[Theorem~3 and Table~I]{Kolchinsky2026}.
An outer localization theorem, however, does not establish attainability
of every point in a proposed region.

Here we address this distinction for four states. We exhibit a continuum
of points in the proposed compressed region that no admissible generator
can realize. These points approach the spectrum of a uniform biased square.
At that endpoint the known localization theorem is rigid. Nearby,
the triangle affinities created by adding chords restrict the quadratic
spectral response. Controlling these new constraints over all supports,
rather than assuming a fixed ring, is essential to the exclusion.

\subsection{Conventions and the spectral-set question}

Use a row generator $Q=(q_{ij})_{i,j=0}^3$: for $i\ne j$, $q_{ij}\geq0$
is the rate from $i$ to $j$, and $q_{ii}=-\sum_{j\ne i}q_{ij}$.
Column probabilities obey $\dot p=Q\trans p$. Transposition preserves
the spectrum, and backward observables satisfy $\dot f=Qf$.
The maximum escape rate is
\begin{equation}
 r(Q)=\max_i\sum_{j\ne i}q_{ij}.
\end{equation}
Rates and eigenvalues have inverse-time units. We normalize by the allowed
maximum escape rate, set that bound to one, and henceforth use
dimensionless rates and eigenvalues. Physical units can be restored by
multiplying every rate and eigenvalue by the same positive rate scale.

For an irreducible generator with paired support
($q_{ij}>0$ if and only if $q_{ji}>0$), define, for every present simple
oriented cycle $\gamma$ of length $\ell_\gamma\geq3$,
\begin{equation}\label{eq:affinity}
 F_\gamma=\log\frac{\prod_{(i\to j)\in\gamma}q_{ij}}
                          {\prod_{(i\to j)\in\gamma}q_{ji}},
 \qquad B(Q)=\max_\gamma\frac{|F_\gamma|}{2\ell_\gamma}.
\end{equation}
Set $B(Q)=0$ when no such cycles occur. On four states their lengths
are three or four; backtracking two-cycles have zero affinity.
Both orientations are included. Stationary-probability factors telescope
around a cycle, so under local detailed balance $F_\gamma$ is the entropy
change in the environment per completed cycle, in units of $k_{\rm B}$.
The spectral argument uses the rate ratios and does not require an
additional reservoir model.

Let $\Sigma_{4,B}$ be the union of individual eigenvalues over this
irreducible, paired-support class with $r(Q)=1$ and $B(Q)=B$.
Let $\Omega'_4$ be the corresponding union over all four-state generators
with $r(Q)=1$, without an affinity restriction and allowing irreversible
or reducible matrices. Define
\begin{equation}
 D_B(x+\ii y)=x+\ii\tanh(B)y.
\end{equation}
The equality under consideration is
\begin{equation}\label{eq:question}
 \Sigma_{4,B}\stackrel{?}{=}D_B(\Omega'_4)
 \qquad\text{for every finite }B\geq0.
\end{equation}
This is a question about individual eigenvalues, not prescribed complete
spectra. We keep exact maximum rate and exact affinity on the left of
\eqref{eq:question}. The exclusion below holds in a larger bounded class,
so it also covers these exact constraints.

The unrestricted spectral region is related to the stochastic
Karpelevich region by uniformization. We need only one explicit
zero-diagonal stochastic companion to prove membership in
$\Omega'_4$. We therefore do not assume a global equivalence between
the exact-rate and bounded-rate spectral unions.

\section{The missing arc}\label{sec:result}

Fix $\Bstar=(\log3)/2$, for which $\tanh\Bstar=1/2$.
Define $\K$ as the set of four-state row generators with $r(Q)\leq1$
and the polynomial inequalities
\begin{equation}\label{eq:cycleproducts}
 \prod_{(i\to j)\in\gamma}q_{ij}
 \leq 3^{\ell_\gamma}\prod_{(i\to j)\in\gamma}q_{ji}
 \quad\text{for every oriented simple three- or four-cycle }\gamma.
\end{equation}
These conditions are imposed even when a product vanishes. The set $\K$
is a closed subset of the bounded rate cube, hence compact. It contains
the closure of the paired-support class with $r\leq1$ and $B\leq\Bstar$;
we do not need or claim equality with that closure. One-way edges between
components are allowed, but no positive directed cycle can contain an
irreversible edge: its positive forward product in
\eqref{eq:cycleproducts} forces every reverse rate positive.

\begin{theorem}\label{thm:missing}
There exists $\eta\in(0,1/2)$ such that every $Q\in\K$ and every
$\lambda\in\spec(Q)$ with $\delta=-1-\Rea\lambda\in(0,\eta)$ obey
\begin{equation}\label{eq:mainbound}
 |\Ima\lambda|\leq\frac12-\frac\delta2+\frac{\delta^2}{4}.
\end{equation}
For every $\delta\in(0,\eta)$, the point
\begin{equation}\label{eq:compressed}
 \Lambda(\delta)=-1-\delta+
       \frac{\ii}{2}\sqrt{1-2\delta+3\delta^2}
\end{equation}
belongs to $D_{\Bstar}(\Omega'_4)$ and lies outside
$\bigcup_{Q\in\K}\spec(Q)$. Each such point has positive distance
from that spectral union. In particular, \eqref{eq:question} is false.
\end{theorem}

The proof gives the stronger uniform asymptotic estimate
\begin{equation}\label{eq:strongbound}
 h\leq\frac{\delta^2}{5}+C_2\delta^3,
 \qquad h=\Ima\lambda-\frac12+\frac\delta2,
\end{equation}
for $h\geq0$ and sufficiently small positive $\delta$, with one finite
constant $C_2\geq0$ over the full class. By comparison,
\begin{equation}\label{eq:arcexpansion}
 \Ima\Lambda(\delta)=\frac12-\frac\delta2+
                  \frac{\delta^2}{2}+O(\delta^3).
\end{equation}
The coefficient $1/5$ is an upper bound; its attainability is not proved.
Neither $\eta$ nor the positive distance in the theorem is supplied
numerically. The argument does not decide an arbitrary fixed interior
point of the arc.

\section{Rigidity of the limiting generator}\label{sec:rigidity}

We use the following published localization input
\cite[Theorem~3]{Kolchinsky2026}. For a nonreal eigenvalue $x+\ii y$,
$y>0$, of a finite generator with maximum escape rate $r$, and each
$\tau\in[0,1]$, a certifying cycle $\gamma$ of length $\ell$ and
eigenvector winding $w\in\{1,\ldots,\lceil\ell/2\rceil-1\}$ satisfies
\begin{equation}\label{eq:sourcebound}
 \tau\frac{y}{-x}\tan\frac{\pi w}{\ell}
 +(1-\tau)\frac{y}{2r+x}\cot\frac{\pi w}{\ell}
 \leq\tanh\frac{F_\gamma}{2\ell}.
\end{equation}
The theorem applies to arbitrary finite generators, and its
left-eigenvector convention for $Q\trans$ is the right-eigenvector
equation $Qf=\lambda f$ used here. For $Q\in\K$ every positive cycle
has finite affinity by \eqref{eq:cycleproducts}.

For a nonreal eigenvalue, $-2r<x<0$. Choose $\tau=-x/(2r)$ in
\eqref{eq:sourcebound}. Its left side becomes
$y/[r\sin(2\pi w/\ell)]$. Thus
\begin{equation}\label{eq:frequencybound}
 y\leq r/2\leq1/2 \qquad (Q\in\K).
\end{equation}
If $\lambda_0=-1+\ii/2$ occurs, necessarily $r=1$. With $\tau=1/2$,
the left side of \eqref{eq:sourcebound} for a three-cycle is
$1/\sqrt3>1/2$. A certifying cycle must therefore have length four,
winding one, and affinity $4\log3$.

\begin{lemma}\label{lem:rigidity}
Every generator in $\K$ with eigenvalue $\lambda_0=-1+\ii/2$ is, up
to a permutation of the states, the uniform biased square
\begin{equation}\label{eq:ring}
 Q_0=\begin{pmatrix}
 -1&3/4&0&1/4\\
 1/4&-1&3/4&0\\
 0&1/4&-1&3/4\\
 3/4&0&1/4&-1
 \end{pmatrix}.
\end{equation}
\end{lemma}

\begin{proof}
We retain the equality information in the finite-affinity proof of
\eqref{eq:sourcebound} \cite[SM3, Eq.~(S55), and SM5]{Kolchinsky2026}.
On the certifying Hamiltonian cycle, write the eigenvector ratios as
$f_{j+1}/f_j=\exp(-\xi_j+\ii\phi_j)$. At the present endpoint the
source parameters are $\gamma=0$ and $\sigma=1/2$, and its cycle
construction yields
\begin{equation}\label{eq:jensen}
 4\log3\geq F_\gamma\geq\sum_{j=0}^3 J(\phi_j,\xi_j)
 \geq 4J(\pi/2,0)=4\log3,
\end{equation}
where
\begin{equation}\label{eq:J}
 J(\phi,\xi)=\log\frac{e^\xi\sin\phi+1/2}
                              {e^{-\xi}\sin\phi-1/2},
 \qquad 0<\phi<\pi,\quad \xi<\log(2\sin\phi).
\end{equation}
The phase and log-modulus averages are $(\pi/2,0)$.
The open domain is convex, and $J$ is strictly convex there; an explicit
Hessian check is given in Appendix~\ref{app:jensen}.
Equality in \eqref{eq:jensen} forces $(\phi_j,\xi_j)=(\pi/2,0)$
for all four edges. Thus $f$ is a nonzero multiple of
$(1,\ii,-1,-\ii)\trans$ in cycle order.

Let $a_j=q_{j,j+1}$, $b_j=q_{j,j-1}$ and $c_j=q_{j,j+2}$, with
indices modulo four. The certifying ring edges and their reverses are
positive. Dividing the $j$th eigenvector equation by $f_j$ gives
\begin{equation}\label{eq:endpointeigen}
 a_j+b_j+2c_j=1,\qquad a_j-b_j=1/2.
\end{equation}
Hence $a_j=3/4-c_j$, $b_j=1/4-c_j>0$, and
\begin{equation}
 \frac{a_j}{b_j}=\frac{3/4-c_j}{1/4-c_j}\geq3,
\end{equation}
with strict inequality for $c_j>0$. The Hamiltonian cycle has
$\prod_j a_j/b_j\leq3^4$. Every chord rate therefore vanishes, giving
\eqref{eq:ring}. This also excludes other reducible or zero-rate
limits at the endpoint.
\end{proof}

The characteristic polynomial of $Q_0$ is
$\lambda(\lambda+2)((\lambda+1)^2+1/4)$. Its spectrum reproduces the
uniform-cycle formula in Ref.~\cite[Eq.~(S66)]{Kolchinsky2026}; in
particular, $\lambda_0$ is simple. Compactness now implies that any
sequence of generators in $\K$ with eigenvalues tending to $\lambda_0$
has distance tending to zero from the finite set of labelled copies of
$Q_0$.

\section{Uniform expansion with escape and affinity slack}\label{sec:expansion}

In a fixed neighborhood of $Q_0$ all eight ring rates remain positive.
Let $\av x=\tfrac14\sum_{j=0}^3x_j$, and introduce
\begin{align}
 s_j&=1-\sum_{k\ne j}q_{jk}\geq0,&
 v_j&=\tfrac12\log(a_j/b_j)-\Bstar,\nonumber\\
 m&=-\av v\geq0,&w_j&=v_j+m,&\av w&=0.\label{eq:coordinates}
\end{align}
Here $m\geq0$ follows exactly from the Hamiltonian cycle inequality.
The four nonnegative $c_j$ describe both directions of both chords.
All nearby generators are represented by
\begin{align}
 a_j&=\frac{1-s_j-c_j}{2}\bigl[1+\tanh(\Bstar+w_j-m)\bigr],\nonumber\\
 b_j&=\frac{1-s_j-c_j}{2}\bigl[1-\tanh(\Bstar+w_j-m)\bigr].\label{eq:logistic}
\end{align}
Set $a=3/4$, $b=1/4$, and $u_j=2ab\,w_j=3w_j/8$, so $\av u=0$.
With fixed Euclidean norms, write
\begin{equation}
 \Nloc=\|c\|+\|s\|+\|u\|+m.
\end{equation}
All remainder constants below are uniform in a sufficiently small fixed
neighborhood. Indeed, \eqref{eq:logistic} extends analytically to signed
coordinates in an open neighborhood. Simplicity of $\lambda_0$ gives one
analytic upper eigenvalue branch there, with bounded third derivatives
on a smaller closed ball. Feasibility is imposed after this analytic
extension; no differentiable parametrization of feasible matrices is
assumed.

The homogeneous linear coefficient of that eigenvalue is
\begin{equation}\label{eq:linear}
 \lambda^{(1)}=-(1+\ii/2)\av c+(1-\ii/2)\av s-\frac{3\ii}{4}m.
\end{equation}
Consequently, for $\delta=-1-\Rea\lambda$ and
$h=\Ima\lambda-1/2+\delta/2$,
\begin{align}
 \delta&=\av c-\av s+O(\Nloc^2),\label{eq:deltaexp}\\
 h&=-\av s-\frac34m+K(c,u)
       +O\bigl((\av s+m)\Nloc+\Nloc^3\bigr).\label{eq:hexp}
\end{align}
Every quadratic term omitted from $K$ contains an escape or affinity
slack. Nonnegativity of $s_j$ and $m$ bounds their sum by a constant
times $(\av s+m)\Nloc$. Thus \eqref{eq:hexp} retains unequal escape
rates and unused affinity explicitly.

\subsection{The quadratic form}

The following formulas specify the complete slack-free quadratic form.
Let $V$ have zero diagonal and entries
\begin{equation}\label{eq:V}
 V_{j,j+1}=-ac_j+u_j,\qquad
 V_{j,j-1}=-bc_j-u_j,\qquad V_{j,j+2}=c_j.
\end{equation}
Let $U_{jk}=\ii^{jk}/2$, $M=U^*VU$, and
$(\ell_0,\ell_1,\ell_2,\ell_3)=(0,-1+\ii/2,-2,-1-\ii/2)$.
The standard simple-root expansion in this diagonal basis gives
\begin{align}
 W&=\sum_{k\ne1}\frac{M_{1k}M_{k1}}{\ell_1-\ell_k},\nonumber\\
 K(c,u)&=\Ima W-\tfrac12\Rea W
       -\frac14\sum_{j=0}^3\left(2c_ju_j+\frac83u_j^2\right).
       \label{eq:quadratic}
\end{align}
The last term comes from the nonlinear coordinates, not from the
resolvent: $\tanh(\Bstar+v)=1/2+3v/4-3v^2/8+O(v^3)$, and its
product with $1-s_j-c_j$ gives the intrinsic quadratic rate correction.
Appendix~\ref{app:invariants} supplies a second representation of
\eqref{eq:quadratic} using characteristic invariants.

Eliminate $u_3=-u_0-u_1-u_2$ and write
$\widehat u=(u_0,u_1,u_2)\trans$. Direct expansion of
\eqref{eq:quadratic} gives the Hessian
\begin{equation}\label{eq:hessian}
 H=\begin{pmatrix}
 -14/3&-4/3&-10/3\\
 -4/3&-8/3&-4/3\\
 -10/3&-4/3&-14/3
 \end{pmatrix}.
\end{equation}
Its leading principal minors are $-14/3$, $32/3$, and $-640/27$;
therefore it is negative definite. The unique maximizer at fixed $c$ is
\begin{align}
 u_0^*&=(-3c_0+6c_1+3c_2-6c_3)/40,\nonumber\\
 u_1^*&=(-6c_0-3c_1+6c_2+3c_3)/40,\label{eq:optimizer}\\
 u_2^*&=-u_0^*,\qquad u_3^*=-u_1^*.\nonumber
\end{align}
Define opposite-pair sums and differences by
\begin{equation}\label{eq:pairs}
 S=c_0+c_2,\quad T=c_1+c_3,\quad
 D=(c_0-c_2)/2,\quad E=(c_1-c_3)/2.
\end{equation}
Exact completion of the square yields
\begin{align}
 K(c,u)&=K_{\max}(c)+\tfrac12(\widehat u-\widehat u^*)\trans
                         H(\widehat u-\widehat u^*),\label{eq:completion}\\
 K_{\max}(c)&=\frac25(D^2+E^2)-\frac5{64}(S-T)^2.\label{eq:Kmax}
\end{align}
These are homogeneous polynomial identities, with no constraint on
$\sum_jc_j$. In particular, the optimization does not impose the
other cycle inequalities or assert that its maximizing direction is
realized by an exact feasible family.

\subsection{A coercive bound fixes the scale of approach}\label{sec:bootstrap}

A quadratic expansion alone would not constrain arbitrary approaches to
the endpoint: initially a rate deviation could be much larger than
$\delta$. We now control that possibility. Put
\begin{equation}
 C=\av c,\qquad L=\av s+m,\qquad U_0=\|u\|.
\end{equation}
Since $c_j\geq0$, $\|c\|\leq4C$. Negative definiteness in
\eqref{eq:completion}, linearity of $u^*(c)$, and Young's inequality
give constants $A,k>0$ such that
\begin{equation}\label{eq:coerciveK}
 K(c,u)\leq AC^2-kU_0^2.
\end{equation}
Norm equivalence and nonnegative slacks also give
$\Nloc\leq C_0(C+L+U_0)$. In a ball $\Nloc\leq\varepsilon$,
\begin{equation}\label{eq:absorption}
 L\Nloc\leq\varepsilon L,\qquad
 \Nloc^3\leq C_3\varepsilon(C^2+U_0^2)
                   +C_4\varepsilon^2L.
\end{equation}
For example, cube the preceding norm bound, absorb mixed terms into
the individual cubes, and use $C,U_0,L=O(\varepsilon)$.
Choosing $\varepsilon$ sufficiently small in
\eqref{eq:hexp}--\eqref{eq:absorption} absorbs the slack and $U_0^2$
remainders into the negative terms. For fixed $A_1,k_1>0$,
\begin{equation}\label{eq:bootstrapineq}
 h\leq A_1C^2-k_1(L+U_0^2).
\end{equation}

On the part with $h\geq0$, this implies
$L=O(C^2)$, $U_0=O(C)$, and $\Nloc=O(C)$. Equation
\eqref{eq:deltaexp} then gives $\delta=C+O(C^2)$.
If $C=0$, \eqref{eq:bootstrapineq} forces every coordinate to vanish
and hence $\delta=0$. If $C>0$ is sufficiently small, then
$C/2\leq\delta\leq3C/2$. Substitution proves the uniform estimates
\begin{equation}\label{eq:bootstrap}
 C=\delta+O(\delta^2),\qquad
 \Nloc=O(\delta),\qquad L=O(\delta^2).
\end{equation}
No smoothness of a sequence of feasible generators enters this argument.

\section{Triangle constraints and exclusion}\label{sec:exclusion}

\subsection{The triangles opened by chords}

The cycle $j\to j+1\to j+2\to j$ in
\eqref{eq:cycleproducts} gives
\begin{equation}\label{eq:triangle}
 a_ja_{j+1}c_{j+2}\leq27b_{j+1}b_{j+2}c_j.
\end{equation}
The forward triangle starting at $j+2$ supplies the opposite comparison.
Since $a_j=3/4+O(\Nloc)$ and $b_j=1/4+O(\Nloc)$, and ring
rates stay bounded away from zero, these inequalities imply
\begin{equation}\label{eq:chordratios}
 c_{j+2}\leq3[1+O(\Nloc)]c_j,\qquad
 c_j\leq3[1+O(\Nloc)]c_{j+2}.
\end{equation}
Only positive ring rates were divided out. Thus the estimates remain
valid on zero-chord faces; one zero direction forces its reverse to
vanish. In the variables \eqref{eq:pairs},
\begin{equation}\label{eq:pairbound}
 |D|\leq S/4+O(\Nloc)S,\qquad
 |E|\leq T/4+O(\Nloc)T.
\end{equation}
Combining this with the nonpositive square in
\eqref{eq:completion} gives
\begin{align}
 K(c,u)&\leq\frac{(S+T)^2}{80}
       -\frac{21}{320}(S-T)^2+O(\Nloc)C^2\nonumber\\
       &\leq\frac{C^2}{5}+O(\Nloc)C^2,\label{eq:trianglecurvature}
\end{align}
because $S+T=4C$. The other cycle constraints can only decrease
the allowed class.

On $h\geq0$, apply \eqref{eq:bootstrap} to
\eqref{eq:hexp} and \eqref{eq:trianglecurvature}. The negative
linear slack terms may be dropped, while all remaining errors are
$O(\delta^3)$. This proves
\begin{equation}\label{eq:localfinal}
 h\leq\frac{\delta^2}{5}+O(\delta^3)
\end{equation}
uniformly near every labelled copy of $Q_0$.

\subsection{Compactness covers the full class}

Suppose generators in $\K$ with $\delta\to0^+$ and $h\geq0$
could remain outside fixed small neighborhoods of those finitely many
copies of $Q_0$. By \eqref{eq:frequencybound},
$1/2-\delta/2\leq\Ima\lambda\leq1/2$, so the eigenvalues tend to
$\lambda_0$. A compact subsequence of generators converges, and the
limit has eigenvalue $\lambda_0$ by continuity of the characteristic
polynomial. Lemma~\ref{lem:rigidity} puts that limit in the finite
set, a contradiction.

Thus one uniform constant $C_2$ in \eqref{eq:strongbound} covers all
such generators for sufficiently small positive $\delta$. Shrink
$\eta<1/2$ further so that $C_2\eta\leq1/20$.
Then \eqref{eq:mainbound} holds for upper eigenvalues with $h\geq0$.
If $h<0$, it is immediate. Real eigenvalues satisfy the bound as well,
and negative imaginary parts follow by conjugation. This proves
the exclusion bound for the entire compact class, including all
support changes and reducible limits.

\subsection{Exact membership of the compressed comparison arc}

For $0\leq t\leq1$ consider the column-stochastic companion
\begin{equation}\label{eq:companion}
 P(t)=\begin{pmatrix}
 0&0&0&1-t\\
 1&0&0&t\\
 0&1&0&0\\
 0&0&1&0
 \end{pmatrix}.
\end{equation}
Its diagonal vanishes, so $P(t)-I$ is a column generator with every
escape rate exactly one. The row generator $P(t)\trans-I$ therefore
belongs to the unrestricted class defining $\Omega'_4$.
For $0<t<1$ it is irreducible. Its characteristic polynomial before
the shift is
\begin{equation}\label{eq:companionpoly}
 z^4-tz-(1-t)=(z-1)(z^3+z^2+z+1-t).
\end{equation}
This is the four-state arc polynomial obtained from the Karpelevich
parametrization in Ref.~\cite[Supplementary Eq.~(S2)]{Xu2025} for
denominators three and four. The explicit matrix proves the needed
membership without appealing to a global description of $\Omega'_4$.

For $0<\delta<1/2$, set
\begin{equation}\label{eq:tandz}
 t(\delta)=4\delta-8\delta^2+8\delta^3,\qquad
 z(\delta)=-\delta+\ii\sqrt{1-2\delta+3\delta^2}.
\end{equation}
Direct substitution into the cubic in \eqref{eq:companionpoly}
gives zero. Moreover,
$t'(\delta)=4(6\delta^2-4\delta+1)>0$, with $t(0)=0$ and
$t(1/2)=1$. This traces the upper arc from $\ii$ to
$e^{2\pi\ii/3}$. In particular, $z(\delta)-1\in\Omega'_4$ and
$D_{\Bstar}(z(\delta)-1)=\Lambda(\delta)$.

To compare with \eqref{eq:mainbound}, both quantities to be squared
are positive on this interval, and
\begin{equation}\label{eq:separation}
 (1-2\delta+3\delta^2)-(1-\delta+\delta^2/2)^2
 =\delta^2+\delta^3-\delta^4/4>0.
\end{equation}
Therefore $\Ima\Lambda(\delta)>1/2-\delta/2+\delta^2/4$.
For $0<\delta<\eta$ these points are absent from $\spec(\K)$,
where $\spec(\K)$ denotes the union of individual spectra.
Finally,
\begin{equation}
 \{(Q,\lambda):Q\in\K,\ |\lambda|\leq2,
                  \det(\lambda I-Q)=0\}
\end{equation}
is compact and projects onto $\spec(\K)$; the modulus bound follows
from Gershgorin's theorem. Each excluded point has positive distance
from this compact union. Since the exact-rate, exact-affinity class
is contained in $\K$, Theorem~\ref{thm:missing} follows.

\section{Physical meaning and scope}

At the endpoint, the uniform square has clockwise and counterclockwise
rates $3/4$ and $1/4$. Increasing the normalized decay displacement
$\delta=-1-\Rea\lambda$ while staying above the endpoint tangent
requires chord rates of order $\delta$. Those transitions also generate
three-state cycles. Their affinity budgets restrict opposite chord rates
to a ratio approaching the interval $[1/3,3]$, which limits the
second-order frequency response. This explains why a deformation based
only on compressing the unrestricted shape overestimates attainability.

The distinction from earlier results is specific. The perturbation
class in Uhl and Seifert \cite[Sec.~4.1, condition~(i), and
Sec.~4.2]{UhlSeifert2019} retains the existing edges of a ring.
It does not incorporate the triangles opened here. The more recent
kinetic refinement of Xu and Meng \cite[Secs.~II--III and
Corollary~3]{XuMeng2026} assumes a unicyclic generator and incorporates
geometric means of forward and backward rates. Uniform-cycle rigidity
within that model is not the contribution claimed here. The general
localization theorem \cite{Kolchinsky2026} is our starting point;
the additional all-support curvature estimate gives the nonattainment
conclusion. These comparisons concern the stated results and assumptions
of the cited works, without claiming an exhaustive classification of
related spectral inequalities.

The theorem concerns an exact, time-independent four-state Markov model.
Second-order perturbation theory is part of a proof with a uniform
remainder, not a finite-order replacement of the dynamics. It excludes
normalized decay--frequency pairs even from the closure of the
bounded-affinity class. It does not give the full attainable boundary,
establish $1/5$ as its sharp curvature, or provide an effective value of
$\eta$. In particular, the point at the companion parameter $t=1/2$
is undecided by this proof. No general inclusion in the compressed
Karpelevich region is proved or disproved: the set equality fails by
nonattainment. Determining a quantitative validity interval or the actual
sharp local boundary remains a separate problem.

\section*{Acknowledgements and generative AI use}

Substantial generative-AI assistance through OpenAI Codex was used for
problem exploration, analytic reasoning, symbolic and exact-arithmetic
code, source checks, internal consistency checks, and preparation of this
draft. The writing environment identifies its model as GPT-6. Exact
model/version information for earlier automated stages remains to be
confirmed. Internal checks in separate AI contexts do not constitute
external peer review. Human scientific and bibliographic review of this
local draft has not yet been completed.

\section*{Data and code availability}

The accompanying reproducibility material contains the exact symbolic
and rational-arithmetic checks, their expected outputs, and instructions
for regeneration. No experimental data or numerical parameter scan is
used in the proof. The scripts check algebraic identities; the uniform
remainder, scale control, and compactness arguments are given in the text.

\appendix
\section{Strictness of the endpoint Jensen step}\label{app:jensen}

For completeness, the strict convexity needed in
Lemma~\ref{lem:rigidity} can be checked explicitly from \eqref{eq:J}.
Put $v=\sin\phi$, $\rho=e^\xi$, $N=\rho v+1/2$,
and $M=v/\rho-1/2$. On the open domain in \eqref{eq:J},
$0<v\leq1$, $0<\rho<2v$, and $N,M>0$. Direct differentiation gives
\begin{equation}
 \partial_\xi^2J=\frac{\rho v}{2N^2}+
                         \frac{v}{2\rho M^2}>0
\end{equation}
and the specialization of the Hessian identity of
Ref.~\cite[SM5]{Kolchinsky2026},
\begin{equation}\label{eq:detJ}
 \det\nabla^2J=
 \frac{\frac54v^2\bigl((\rho+\rho^{-1})/2\bigr)^2
 +\frac14(1-v^2)\bigl((\rho-\rho^{-1})/2-2v\bigr)^2}
 {N^3M^3}>0.
\end{equation}
The first numerator term is strictly positive everywhere in the domain.
The Hessian is positive definite, and the domain is convex because
$\log(2\sin\phi)$ is concave on $(0,\pi)$. Hence equality in Jensen's
inequality forces identical phase and log-modulus pairs.

\section{Reconstruction from characteristic invariants}\label{app:invariants}

The quadratic form may be verified without an eigenbasis. With $s=m=0$,
write $A=Q+I$. Its diagonal is zero and its row sums are one. Let
\begin{equation}
 P_2=\sum_{i<j}A_{ij}A_{ji},\qquad
 P_3=\sum_{i<j<k}\bigl(A_{ij}A_{jk}A_{ki}+A_{ik}A_{kj}A_{ji}\bigr).
\end{equation}
The trace identities give the $z^2$ and $z$ coefficients of
$\det(zI-A)$ as $-P_2$ and $-P_3$. Since one is an eigenvalue,
factoring $z-1$ leaves
\begin{equation}\label{eq:invariantcubic}
 f(z)=z^3+z^2+(1-P_2)z+1-P_2-P_3.
\end{equation}

Introduce a formal scalar $\epsilon$ by replacing $(c,u)$ with
$(\epsilon c,\epsilon u)$ in \eqref{eq:logistic}. Through degree two,
\begin{align}
 A_{j,j+1}&=\tfrac34+\epsilon(u_j-3c_j/4)
        -\epsilon^2(c_ju_j+4u_j^2/3)+O(\epsilon^3),\nonumber\\
 A_{j,j-1}&=\tfrac14+\epsilon(-u_j-c_j/4)
        +\epsilon^2(c_ju_j+4u_j^2/3)+O(\epsilon^3),\label{eq:intrinsic}\\
 A_{j,j+2}&=\epsilon c_j.\nonumber
\end{align}
Write $f=f_0+\epsilon f_1+\epsilon^2f_2+O(\epsilon^3)$ and
$z=z_0+\epsilon z_1+\epsilon^2z_2+O(\epsilon^3)$, where $z_0=\ii/2$.
Since $f_0'(z_0)\ne0$,
\begin{equation}\label{eq:implicit}
 z_1=-\frac{f_1(z_0)}{f_0'(z_0)},\qquad
 z_2=-\frac{f_2(z_0)+f_1'(z_0)z_1+\tfrac12f_0''(z_0)z_1^2}
                 {f_0'(z_0)}.
\end{equation}
Substitution gives $z_1=-(1+\ii/2)\av c$ and
$K(c,u)=\Ima z_2-\Rea z_2/2$, reproducing
\eqref{eq:hessian}--\eqref{eq:Kmax} as identities in all seven
independent variables $(c_0,c_1,c_2,c_3,u_0,u_1,u_2)$.

The supplied symbolic program implements these invariants. A separate
program uses Gaussian rational arithmetic and a bordered linear system
for the eigenvector derivative. Because the resulting coefficient is a
homogeneous quadratic polynomial in seven variables, the seven basis
directions and their twenty-one pair sums determine all twenty-eight
coefficients exactly. Five further first-order directions check the four
escape slacks and the mean-affinity slack in \eqref{eq:linear}.
This is a complete check of the finite quadratic identity, not a fit or
a numerical argument for the universal validity interval.

\bibliographystyle{unsrtnat}
\bibliography{references}
\end{document}
